\section{Introduction}\label{sec:introduction}

A motion-imitation controller should reproduce an action without assuming
that every character has the same body. Changes in limb lengths, mass
distribution, and collision geometry alter both the reference trajectory and
the dynamics required to follow it: replaying an identical joint sequence
across bodies does not preserve foot placement, clearance, or balance, while
optimizing only spatial agreement can obscure whether the intended
articulation survives. Learning across body shapes therefore requires a
consistent relationship between source motion, shape-specific reference, and
simulated body.

Existing work provides large motion-capture collections, body-shape-conditioned
motion generation, and scalable physics-based control~\cite{amass,humanml3d,humos,deepmimic,phc,protomotions}.
We build on these components to study a practical question: how can one
policy imitate a large collection of shape-specific motions while preserving
traceable data identities and a meaningful evaluation protocol? Our
contribution is this integration and its empirical assessment, not a claim
that morphology conditioning or attention is new.

Our data pipeline extends AMASS-derived HumanML3D clips to 128 body
configurations. HUMOS generates the corresponding reference motions; SMPLSim
constructs simulation assets from the same body parameters. A conservative
refinement pass smooths rotations, reduces stance sliding through root
translation, and enforces clearance of the modeled collision geometry. The
inherited descriptions connect each source action to its generated variants,
supporting paired analysis and future language-based uses, although the
controller studied here consumes reference motion, not text.

The selected controller combines current state, dilated state history, and
future reference targets through temporal slot/type attention, with body
parameters and a previous action concatenated into separate actor and critic
heads. The main experiment trains directly on multiple shapes; neutral-body
pretraining and mixed-source reward training were explored during
development but are not part of the final method. A 150-clip development
subset makes design exploration affordable and precedes full-scale training
on a versioned split.

This paper makes four contributions:
\begin{enumerate}
    \item A derived, text-associated multi-morphology corpus, with an explicit
    construction procedure linking source clips, generated references, body
    assets, and reproducible splits. Its captions are inherited rather than
    newly annotated for each body.
    \item A single morphology-conditioned imitation system operating across
    18,855 training clips and 128 body configurations, integrating temporal
    observations, matching simulation assets, reference refinement, and
    memory-bounded motion streaming.
    \item An empirical protocol separating architecture, reference quality,
    and evaluation correctness: an eight-run architecture/reference grid on a
    150-clip development set, a full-scale validation and independent test
    evaluation of the selected checkpoint, and a caption-grounded analysis of
    the residual failures.
    \item A physical analysis of how body shape and motion type jointly set
    control demand, replaying the trained policy across every clip/shape pair
    of the development corpus and testing -- rather than assuming -- whether
    dynamic motions are more shape-sensitive than quasi-static ones.
\end{enumerate}

The selected checkpoint reaches 98.47\% success on validation and 99.14\% on
the untouched test split under the protocol in
Section~\ref{sec:evaluation_protocol}. These results concern held-out clip
identities on trained body configurations, with one body evaluated per clip;
they do not establish success for every shape or for unseen bodies.
Section~\ref{sec:physics} tests the morphology$\times$motion-class
interaction directly and finds no statistically detectable effect for this
checkpoint, once a normalization artifact in the naive comparison is
corrected. An evaluation on body shapes drawn from the trained coefficient
range but excluded from the 128 training configurations
(Section~\ref{sec:unseen_morphology}) remains future work.
